% ============================================================================
%  CAPÍTULO V: CONCLUSIONES
% ============================================================================

\section{Conclusiones}

La implementación del servidor de correo electrónico y la planificación del
servidor de archivos permitieron comprender de manera integral el
funcionamiento de los servicios de red en una organización. A continuación se
presentan las conclusiones más relevantes.

\subsection{Resultados obtenidos}

El servidor de correo fue configurado exitosamente con las siguientes
características verificadas:

\begin{itemize}
  \item Servicio SMTP habilitado en los puertos 25 y 587, con autenticación
        obligatoria y protección STARTTLS en el puerto de submission.
  \item Servicios IMAP (143) y POP3 (110) operativos, permitiendo a los
        clientes acceder a sus buzones de correo de forma remota.
  \item Resolución de nombres funcional para el dominio \texttt{sudoers.lan},
        incluyendo registro MX y registros A para los servicios de correo.
  \item Usuarios creados y autenticados correctamente a través del subsistema
        SASL integrado en Dovecot.
  \item Envío y recepción de correo entre usuarios del dominio verificado
        satisfactoriamente.
  \item Servidor de archivos Samba (SMB) configurado con carpetas compartidas
        y permisos por usuario (desarrollado por otro integrante del equipo).
\end{itemize}

\subsection{Dificultades encontradas}

Durante la configuración se encontraron las siguientes dificultades:

\begin{itemize}
  \item \textbf{Certificados auto-firmados:} el uso de certificados SSL
        auto-firmados genera avisos de seguridad en los clientes de correo,
        requiriendo la aceptación manual de excepciones. En un entorno de
        producción se recomienda utilizar certificados emitidos por una
        autoridad de certificación confiable.
  \item \textbf{Orden de inicio de servicios:} el servidor de correo depende
        del correcto funcionamiento del servidor DNS para resolver registros
        internos, lo que requiere un orden específico en el inicio de los
        servicios del sistema.
  \item \textbf{Coincidencia de configuración Postfix-Dovecot:} ambos
        servicios deben coincidir en el formato de buzón (Maildir) y en la
        ruta del socket de autenticación SASL, lo que requiere atención
        detallada durante la configuración.
\end{itemize}

\subsection{Aprendizajes}

El taller permitió desarrollar competencias prácticas en administración de
servicios de red, comprendiendo la arquitectura cliente-servidor del correo
electrónico y la importancia de la configuración correcta de los registros
DNS para la operación de servicios internos. Asimismo, se valoró la
instalación y configuración manual de servicios en Debian, lo que brinda un
entendimiento profundo del funcionamiento de cada componente del sistema
operativo y de la interacción entre los distintos servicios de red.
